
\needspace{2\baselineskip}
\color{uniblue}{\subsection{Блокировка при неисправности цепей напряжения до ввода (БННвв)}}
\color{black}

\begin{enumerate}

\item
Блокировка при неисправности цепей напряжения используется для блокировки органов защит, которые могут работать неправильно при нарушениях в цепях напряжения. В устройстве блокировка при неисправности цепей напряжения ТН до ввода представлена функцией «БННвв».

\item
Функция включает в себя четыре различных алгоритма:
\begin{itemize}
    \item контроль напряжения небаланса U БНН (dU0), полученного из определенной комбинации мгновенных значений фазных напряжений «звезды» и напряжения $3U_0$ «разомкнутого треугольника»;
    \item контроль расчетных значений напряжения и тока обратной последовательности (U2/I2);
    \item контроль снижения линейных напряжений и наличие фазных токов (minU3p);
    \item контроль третьей гармоники напряжения нулевой последовательности (3U03f).
\end{itemize}

\item
Ввод в работу ИО небаланса по напряжению, ИО напряжения обратной последовательности, ИО минимального напряжения и ИО 3-й гармоники напряжения НП осуществляются уставками SGF2, SGF3, SGF4 соответственно.

\item
Срабатывание функции с режимом работы контроля небаланса (уставка SGF2) происходит при возникновении неисправности во вторичных цепях ТН, при этом баланс напряжений обмоток «звезды» и «треугольника» нарушается. Алгоритм рассчитывает среднее значение суммы фазных напряжений за полпериода и сравнивает его со средним измеренным значением $3U_0$ для того же полупериода. Если разность двух параметров превышает заданную уставку, функция срабатывает и активируется соответствующий выходной сигнал «БННвв:~Пуск dU0». Время срабатывания БНН при выявлении небаланса, устанавливается фиксированной выдержкой времени Т6 которая составляет 20 мс.

\item
Напряжение небаланса рассчитывается по формуле:
% ФОРМУЛА
\begin{equation*}
    dU_0=\lvert 3U_{0\textsl{вычисл.}} - \frac{3U_{0\textsl{вычисл.ном.}}}{3U_{0\textsl{изм.ном.}}}3U_{0\textsl{изм.}}  \rvert,
\end{equation*}
\begin{eqexpl}[25mm]
    \item{$3U_{0\textsl{вычисл.}}$} сумма фазных напряжений «звезды»;
    \item{$3U_{0\textsl{изм.}}$} измеренное напряжение $3U_0$ «разомкнутого треугольника»;
    \item{$3U_{0\textsl{вычисл.ном.}}$} номинальное напряжение основной обмотки ТН;
    \item{$3U_{0\textsl{изм.ном.}}$} номинальное напряжение $3U_0$ обмотки «разомкнутого треугольника» ТН.
\end{eqexpl}

\item
Симметричность вторичных напряжений, подводимых к устройству, определяется уровнем напряжения обратной последовательности $U_2$. Для предотвращения ложного блокирования защит и функций автоматики при несимметричных КЗ в первичной сети контролируется уровень тока обратной последовательности $I_2$. Данный алгоритм контролируется уставкой SGF3 позволяя выявлять все виды несимметричных повреждений в цепях ТН. Срабатывание ИО данного алгоритма активирует выходной сигнал «БННвв: Пуск U2I2». Время срабатывания БНН при выявлении междуфазных замыканий в цепях, устанавливается фиксированной выдержкой времени Т2 которая составляет 20 мс.

\item
Если максимальное из трех линейных напряжений ниже уставки с введенным контролем уставкой SGF4, минимальный из трех фазных токов выше уставки, отсутствуют сигналы срабатывания измерительного органа ОПО функций ДЗ либо общего пуска МТЗ по истечении времени нерегулируемой выдержки времени Т3 активируется выходной сигнал «БННвв: Пуск minU3p». Время срабатывания БНН при симметричных замыканиях и симметричных обрывов в цепях «звезды» устанавливается фиксированной выдержкой времени Т3 которая составляет 20 мс.

\item
Уставкой SGF5 осуществляется ввод в работу контроля утроенного напряжения НП третьей гармоники который отслеживает обрыв в цепях «разомкнутого треугольника». Срабатывание ИО информируется сигналом ИО «3U03f», реагирующим на снижение уровня напряжения третьей гармоники в цепях «разомкнутого треугольника» ниже уставки. В данном случае учитывается тот факт, что даже в обычном режиме работы из-за нелинейности ТН в схемах «разомкнутого треугольника» наблюдается заметный уровень третьей гармоники, составляющий около 0,5 В. При обрыве, уровень третьей гармоники значительно снижается, теоретически приближаясь к нулю. Для исключения ложного срабатывания в нормальных и аварийных режимах, не связанных с неисправностью в цепях «разомкнутого треугольника» выполнен контроль уровня минимального напряжения и контроль напряжения ОП и по истечении времени нерегулируемой выдержки времени таймера Т4 формируется выходной сигнал пуска «БННш: Пуск 3U03f». При наличии пусковых условий данного алгоритма, по истечении выдержки времени таймера Т1 формируется сигнал неисправности в цепях напряжения ТН «БННш: Неисправность ЦН». При выявлении неисправностей в цепях «разомкнутого треугольника» формируется блокирующий сигнал для дистанционной защиты.

\item
В устройстве реализован контроль состояния автоматических выключателей ТН до ввода при помощи соответствующих блок-контактов. Сигнал с блок-контактов АВ ТН до ввода заводится на дискретные входы устройства ««Автомат ТНвв 1», «Автомат ТНвв 2». Для предотвращения кратковременных несимметричных режимов,  возникающих  при одновременном  замыкании  силовых  контактов  автомата,  предусмотрена задержка  на  снятие  блокировки  ступеней  защиты  при  включении  автомата  ТН  длительностью  120  мс. Наличие  аварийного  или  принудительного  отключения  АВ  ТН  до  ввода  формируется  выходной  сигнал «БННвв: Автомат ТН откл».

\item
При  наличии  активного  дискретного  сигнала  «Неиспр.  ТНвв»  и  введенной  функции  в  работу формируется сигнал общего срабатывания функции БНН «БННвв: Срабатывание». 

\item
При наличии условий срабатывания функции «БННвв» активируется таймер независимой выдержки времени Т1. По истечении выдержки времени таймера Т1 активируется выходной сигнал «БННвв: Неисправность ЦН».

\item
Ввод функции в работу на этапе параметрирования устройства осуществляется программным переключателем SGF1 «Ввод\_функции». Информация о введенной ступени отображается выходным сигналом «БННвв: Ввод» на выходе функции.

\item
Функция может быть выведена из работы оперативно путем активации сигнала «ОВ БННвв», что характеризуется активным состоянием сигнала «БННвв: Оперативный вывод».

\item
Алгоритм работы функции выполнен в соответствии с рисунком \ref{bnnvv:img1}. В таблице \ref{bnnvv:tbl1} приведены параметры, необходимые для настройки функции.

\vspace{1mm}
\begin{figure}[H]
\centering
\includegraphics[width=1\textwidth,height=1\textheight,keepaspectratio]{\fbpath/304. БННвв/_latex/img/bnnvv.pdf}
\captionof{figure}{Функциональная схема логики БННвв}\label{bnnvv:img1}
\end{figure}

\end{enumerate}
